% ECONOMICS_PROBLEM: {"id": "O17-374", "title": "Formalization without destroying livelihoods", "jel_code": "O17", "source_id": "OP-0376"}
% !TeX program = xelatex
\documentclass[11pt,a4paper]{article}
\usepackage[margin=23mm]{geometry}
\usepackage{fontspec}
\setmainfont{Latin Modern Roman}
\usepackage{amsmath,amssymb,mathtools}
\usepackage{xcolor,enumitem,booktabs,longtable,array,xurl,hyperref}
\definecolor{muted}{HTML}{53636C}
\hypersetup{colorlinks=true,urlcolor=blue,linkcolor=blue}
\setlength{\parindent}{0pt}
\setlength{\parskip}{5pt}
\newcommand{\R}{\mathbb R}
\newcommand{\E}{\mathbb E}
\newcommand{\N}{\mathbb N}
\newcommand{\ind}{\mathbf 1}
\newcommand{\Var}{\operatorname{Var}}
\newcommand{\Cov}{\operatorname{Cov}}
\newcommand{\supp}{\operatorname{supp}}
\DeclareMathOperator*{\argmin}{arg\,min}
\DeclareMathOperator*{\argmax}{arg\,max}
\newcommand{\entrymeta}[3]{{\sffamily\small\color{muted}\textbf{\texttt{#1}}\quad|\quad #2\par #3\par}}
\newcommand{\Problemref}[1]{\texttt{#1}}
\newcommand{\JELref}[2]{\texttt{#2} (JEL \texttt{#1})}
\begin{document}
\section*{Formalization without destroying livelihoods}
\textbf{Problem ID:} \texttt{O17-374}\quad \textbf{JEL:} \texttt{O17}\par
% BEGIN ECONOMICS STATEMENT
\label{problem:OP-0376}
{\sffamily\small\color{muted}Original subject group: Development and poverty\par}
\label{definitions:problem:30007658}
\textbf{Audit classification:} Published research agenda. Named editors, version, and chapter checked. Dynamic firm and worker transitions and policy transition paths are stated gaps. Website citation misspells Gabriel as Gabrielle; editor identity confirms Gabriel.
\subsection*{Definitions and precise research target}
\textbf{Complete editorial problem.} Fix a labor and product market with informal firms and workers. Policy \(p\) specifies registration costs, tax and labor enforcement, access to services, and social-insurance eligibility. Let \(F_{it}(p)\) indicate formal status, \(Y_{it}(p)\) disposable earnings, and \(U_{it}(p)\) a declared welfare measure including work risk and benefits. Let \(E_t(p)\) denote firm entry and exit. Estimate dynamic effects \(E[F_{iT}(p)-F_{iT}(0)]\) and \(E[U_{iT}(p)-U_{iT}(0)]\), including workers and entrepreneurs initially informal and people displaced by enforcement. Determine whether informal activity serves as an entry route toward productive formal businesses and jobs or sustains persistent segmentation. Use linked longitudinal observations and policy variation to separate selection from causal transitions. Registration alone is not the welfare objective: a policy can increase formal shares by destroying marginal livelihoods. Counterfactuals must therefore incorporate exit, wages, prices, migration, and fiscal use of new revenue. Compare feasible policy bundles under explicit budget and enforcement constraints. The cited living review expressly identifies dynamic firm decisions and worker formal--informal transitions as research gaps; this welfare-centered intervention formulation extends those passages without asserting a universally beneficial formalization policy.

\textbf{Definitions and admissible scope.} Informality is multidimensional: record business registration, tax compliance, worker-contract registration, and social-insurance enrollment separately. Define \(F_{it}\) as one declared component or a specified vector. Baseline entrepreneurs and workers remain in the welfare population after exit, job loss, or migration, using explicit follow-up conventions. Earnings include zero after business closure and net out taxes and benefits consistently. A complete policy fixes enforcement probabilities, penalties, available formal-sector services, and fiscal use of new revenue. Increasing the fraction formal by removing informal firms changes the denominator; report counts, hours, income, and welfare alongside formal shares. Transition effects concern the full assignment regime when markets and public resources respond. Fix a baseline cohort, \(T<\infty\), finite-mean real outcomes, and a no-reform comparator \(p=0\). For a vector of formal statuses report each component of \(E[F_T(p)-F_T(0)]\); for a scalar choose one registration criterion beforehand. Use \(U_{iT}(p)\) as a stated utility index or discounted utility sum, with common currency/price and benefit accounting. Entry and exit \(E_t\) is an ordered pair of counts, not their net difference unless so declared. The target is the joint identified set of transition probabilities, income, welfare, and entry/exit responses under feasible policy bundles. A model-based welfare ordering is conditional on its preferences, technology, enforcement, fiscal, and equilibrium-selection primitives.
\subsection*{Variants and assumption changes}
\begin{enumerate}
\item Stepping-stone dynamics (source-motivated): identify age- and asset-specific transition kernels between informal work, formal work, business ownership, and nonemployment, separating true state dependence from persistent heterogeneity.
\item Transition policy and welfare (source-motivated extension): compare registration subsidies with gradual enforcement and social-insurance access. Evaluate transition paths, displacement, and fiscal closure instead of comparing stationary formal shares alone.
\end{enumerate}
\subsection*{References and source audit}
\textbf{Support and access.} Named editors, version, and chapter checked. Dynamic firm and worker transitions and policy transition paths are stated gaps. Website citation misspells Gabriel as Gabrielle; editor identity confirms Gabriel. \textbf{Locator:} VoxDevLit 6(2), April 2025, Chapter 6, paragraphs 3--6 and final paragraph.
\begin{enumerate}\raggedright
\item\hypertarget{definitions:source:30007658:1}{} Gabriel Ulyssea; Matteo Bobba; Lucie Gadenne; Mariaflavia Harari. 2025. \emph{Informality, VoxDevLit 6(2)}. Chapter 6, Final Remarks, paragraphs 3–6.
\url{https://voxdev.org/voxdevlit/informality-issue-2/final-remarks-informality}.
\end{enumerate}

\subsection*{Common conventions: Source mathematical conventions}

These conventions apply to each entry unless explicitly replaced.

\textbf{Models and identification.} A primitive model \(m\) specifies all probability laws, feasible choices, information, equilibrium and measurement rules used by its target. The admissible class \(\mathcal M\) is fixed independently of outcomes. If \(P_O(m)\) is its observable probability law and \(\tau(m)\) the target, its identified set at observable law \(P\) is
\[
 \mathcal I_\tau(P)=\{\tau(m):m\in\mathcal M,\ P_O(m)=P\}.
\]
Point identification means that this set is a singleton. An empty set signals incompatibility of data and assumptions. Coordinate bounds need not describe the joint set. Bounds are sharp only if every retained target is attainable or arbitrarily approachable by compatible models. Finite bounds require explicit restrictions on unobserved quantities; unrestricted confounding, hidden wealth, extrapolation or arbitrary equilibrium selection can yield uninformative sets.

\textbf{Causal interventions.} A potential outcome \(Y_i(p)\) is the outcome under a complete policy regime \(p\), including timing, financing, information and market spillovers. The target population and units are fixed before treatment. With interference, use the complete assignment vector or a justified exposure mapping; individual no-interference notation alone cannot identify an economy-wide intervention. A difference of mean potential outcomes does not require the cross-world joint distribution. Conditioning on post-treatment migration, survival, enrollment or employment changes the estimand and needs separate justification.

\textbf{Statistical statements.} An estimator, confidence set, forecast or test specifies a sampling law, model class and loss/coverage criterion. Uniform \((1-\alpha)\)-coverage means \(\inf_{m\in\mathcal M}\Pr_m\{\tau(m)\in C_n\}\geq1-\alpha\), or an explicitly asymptotic version. The whole parameter space trivially has coverage, so informativeness requires a stated width, risk or power objective. A forecasting improvement is relative to a named benchmark, information set and evaluation population.

\textbf{Equilibrium and optimization.} An equilibrium comprises feasible plans, optimality, beliefs where needed, budget constraints and all market/resource clearing. The primitives or a stated benchmark define each correspondence \(\mathcal E(p;m)\). Nonemptiness is assumed or proved, never inferred just from compact policy sets. A selected-equilibrium target uses a declared selection rule; a robust target quantifies over all permitted equilibria. Use suprema or infima unless attainment is established. An \(\varepsilon\)-optimal policy is feasible and within \(\varepsilon\) of the finite optimum value. Report an empty feasible set separately.

\textbf{Welfare and accounting.} Normative utility, interpersonal normalization, social weights, numeraire, discounting and the use of public receipts are primitives. Behavioral utility need not coincide with the welfare criterion. Household welfare includes ownership income and rents once; adding profits again to compensating variation generally double-counts them. A monetary equivalent is defined by an explicit transfer and price convention, with monotonicity and range conditions. Average effects, mechanism contributions, transfers and welfare losses are different targets. Infinite sums require convergence and dynamic feasible plans require terminal or no-Ponzi conditions.

\textbf{Source versions.} A working-paper date, revision date and journal publication year are distinguished. A DOI for a journal version is not automatically the DOI of an earlier manuscript. Locators refer to the stated version and printed pages where specified. A metadata-only check does not verify a claim about a full-text conclusion. Every entry's source subsection narrows the provenance claim accordingly.
% END ECONOMICS STATEMENT
\end{document}
